% ===========================================================================
%  §3  The Thinking Tax: An Empirical Study
%  analysis_final.tex — NeurIPS 2026
%  Loaded via % ===========================================================================
%  §3  The Thinking Tax: An Empirical Study
%  analysis_final.tex — NeurIPS 2026
%  Loaded via % ===========================================================================
%  §3  The Thinking Tax: An Empirical Study
%  analysis_final.tex — NeurIPS 2026
%  Loaded via % ===========================================================================
%  §3  The Thinking Tax: An Empirical Study
%  analysis_final.tex — NeurIPS 2026
%  Loaded via \input{sections/analysis_final}
% ===========================================================================

Before introducing our method, we present a systematic empirical study of
the costs and failure modes of chain-of-thought reasoning under constrained
token budgets.  All experiments in this section use
\textbf{Qwen3-8B}~\cite{qwen3_2025} with its native thinking mode
(\texttt{enable\_thinking=True}) and \textbf{Qwen3.5-27B} for cross-scale
validation, evaluated on the full GSM8K test set ($n{=}1{,}319$)
unless otherwise noted.\footnote{%
  A 200-sample subset (seed 42) is used for the budget sweep in
  Table~\ref{tab:thinking-tax}; all other analyses use the complete set.}
We control test-time compute via the \texttt{max\_new\_tokens} parameter,
which caps the \emph{total} number of newly generated tokens---including
both the thinking trace (\texttt{<think>...<\!/think>}) and the final answer.
Crucially, thinking mode and non-thinking mode operate under the
\emph{same} total output-token budget~$b$, ensuring a fair comparison:
\texttt{think@$b$} must fit its entire reasoning chain \emph{plus}
its answer within $b$~tokens, while \texttt{nothink@$b$} uses all
$b$~tokens for the answer directly.
Non-thinking results are obtained by setting
\texttt{enable\_thinking=False} in the Qwen3 chat template, which
suppresses the thinking block entirely.
Greedy decoding ($\tau{=}0$) is used throughout.

Our five findings collectively establish that (i)~thinking mode imposes a
severe \emph{tax} at practical token budgets, (ii)~models already emit
reliable signals indicating when this tax is unnecessary, and
(iii)~a large fraction of compute is wasted on questions the model
cannot solve regardless of budget.
These findings motivate the design of \textsc{Town} in
\S\ref{sec:method}.

% ---------------------------------------------------------------------------
\subsection{Finding 1: Non-Thinking Beats Thinking at All Matched Budgets}
\label{sec:finding:tax}
% ---------------------------------------------------------------------------

Our most striking result is that, when the output-token budget is held
constant, \emph{non-thinking mode uniformly dominates thinking mode}
across every budget level we tested.
Table~\ref{tab:thinking-tax} summarizes the comparison.

\input{../results/paper_figures/table_thinking_tax_main}

At budget $b{=}128$, non-thinking achieves \textbf{54.5\%} accuracy while
thinking yields only \textbf{2.0\%}---a 52.5 percentage-point (pp) gap.
At $b{=}256$, the gap narrows but remains enormous:
\textbf{89.0\%} vs.\ \textbf{22.0\%} ($+67.0$\,pp).
Even at $b{=}512$, where thinking has enough room for short chains,
non-thinking still leads by a large margin:
\textbf{94.0\%} vs.\ \textbf{66.5\%} ($+27.5$\,pp).
These results are confirmed on the full GSM8K test set ($n{=}1{,}319$,
$\star$ rows in Table~\ref{tab:thinking-tax}):
\texttt{nothink@256$\star$} achieves \textbf{87.5\%} vs.\ only
\textbf{18.0\%} for \texttt{think@256$\star$}, a
\textbf{69.5\,pp} gap at matched budget.

The root cause is straightforward: a thinking model must first produce a
reasoning chain \emph{inside} its $b$-token budget before emitting an answer.
At low budgets, the chain consumes nearly all available tokens, leaving
insufficient room for a well-formed answer.
A non-thinking model, by contrast, directly generates the solution,
allocating its full budget to answer-bearing content.
This difference manifests in the \emph{Early Stop} column of
Table~\ref{tab:thinking-tax}: at $b{=}256$, 92.0\% of non-thinking
samples terminate naturally before the budget, whereas only 2.0\% of
thinking samples do.
In effect, the non-thinking model has already ``solved'' the problem and
moved on, while the thinking model is still mid-chain.

% Figure 1 (nothink vs thinking) is presented in the Introduction.
% See Figure~\ref{fig:nothink-vs-thinking} for the visual comparison.

We emphasize that this result does \emph{not} imply that chain-of-thought
reasoning is useless---at sufficiently large budgets (e.g.,\ $b{=}4096$),
thinking mode recovers and eventually surpasses non-thinking.
Rather, it shows that under the budget-constrained regime that is
increasingly relevant for cost-sensitive deployment, the overhead of
producing a reasoning chain is a \emph{tax} that must be justified by
the problem's difficulty.

% ---------------------------------------------------------------------------
\subsection{Finding 2: Natural Stop Is a Free Confidence Oracle}
\label{sec:finding:natural-stop}
% ---------------------------------------------------------------------------

When a thinking model generates a complete reasoning chain and emits
a ``Final answer'' marker \emph{before} exhausting its token budget,
we say the sample exhibits a \emph{natural stop}.
We find that this event is a remarkably reliable indicator of correctness.

\input{../results/paper_figures/table_natural_stop_oracle}

At $b{=}512$ with Qwen3-8B, 37.4\% of samples terminate before consuming
95\% of their token budget---we call these \emph{natural stops}
(Table~\ref{tab:natural-stop-oracle}).
Among these 475 samples, accuracy is \textbf{96.3\%}---nearly perfect.
Among samples that exhaust the budget (i.e., are truncated), accuracy is
only 47.5\%, yielding an accuracy gap of \textbf{$+$48.8\,pp}
(Figure~\ref{fig:natural-stop-split}).
A stricter definition---requiring an explicit ``Final answer'' marker
(\texttt{has\_final})---selects only 6.1\% of samples but achieves
93.8\% accuracy, further confirming the signal's reliability.

This phenomenon admits a simple explanation.
A natural stop means the model's \emph{own} generation process
determined that the reasoning chain was complete---the model ``believes''
it has reached a satisfactory answer.
While LLM confidence is notoriously miscalibrated in
general~\cite{chen2024overthinking}, the natural-stop signal is
structurally different from soft confidence scores: it is a
\emph{binary}, \emph{endogenous} event that does not require access to
logits, hidden states, or any post-hoc calibration.
It is, in a precise sense, a \emph{free} confidence oracle---one that
the model provides at zero additional inference cost.

\begin{remark}[Generality]
The natural-stop oracle effect is not limited to Qwen3-8B.
Table~\ref{tab:natural-stop-oracle} shows that DeepSeek-R1-8B exhibits
similar behavior: at $b{=}512$, 85.0\% of GSM8K samples stop naturally,
and on the harder MATH500 benchmark, natural-stop rates increase steadily
with budget (31.2\% $\to$ 55.0\% $\to$ 76.2\% as $b$ scales from 1024
to 4096).
This universality suggests that the phenomenon is an intrinsic property
of reasoning-trained models, not an artifact of a specific architecture.
\end{remark}

\begin{figure}[t]
\centering
\includegraphics[width=0.72\textwidth]{fig9_acc_split_natural_vs_hit}
\caption{%
  \textbf{Natural stop is a free confidence oracle.}
  Accuracy decomposition by stopping behavior for Qwen3-8B on GSM8K.
  Samples that stop naturally before the budget (blue) achieve dramatically
  higher accuracy than those that exhaust the budget (red), across all
  budget levels.  At $b{=}512$, natural-stop samples reach 96.3\% accuracy
  vs.\ 47.5\% for truncated samples---a 48.8\,pp gap.
}
\label{fig:natural-stop-split}
\end{figure}

% ---------------------------------------------------------------------------
\subsection{Finding 3: Token Utilization Collapses at High Budgets}
\label{sec:finding:utilization}
% ---------------------------------------------------------------------------

A complementary perspective on the natural-stop phenomenon comes from
examining \emph{token utilization}: the fraction of the allocated budget
that the model actually consumes.

\input{../results/paper_figures/table_token_utilization}

Table~\ref{tab:token-utilization} reveals a striking pattern.
At $b{=}128$, Qwen3-8B uses 100\% of its budget---every token is
consumed, and the model is almost certainly still mid-reasoning when
generation is forcibly terminated.
At $b{=}256$, utilization remains near-perfect (99.7\%), with only 1.4\%
of samples stopping naturally.
But at $b{=}512$, utilization drops to \textbf{89.8\%}, with an average
of only 460 tokens consumed out of 512 available.
The model is already ``paying'' for $\sim$52 tokens it does not need.

The pattern is even more dramatic for DeepSeek-R1 on GSM8K: utilization
plummets from 103.0\% at $b{=}256$ (the model slightly exceeds its
budget due to token-counting granularity) to \textbf{38.3\%} at
$b{=}1{,}024$, with 97.5\% of samples stopping naturally.
This means the model uses fewer than 400 tokens on average when allocated
1{,}024---more than 60\% of the budget is wasted.

This utilization collapse is a direct consequence of the heterogeneous
difficulty distribution in typical benchmarks: many problems are easy
enough that a short reasoning chain suffices, and additional budget
yields no benefit.
It provides strong empirical motivation for adaptive budget allocation:
rather than setting a single high budget and accepting low utilization,
one should match the budget to the problem's actual difficulty.

% ---------------------------------------------------------------------------
\subsection{Finding 4: The Thinking Tax Scales with Model Size}
\label{sec:finding:model-size}
% ---------------------------------------------------------------------------

One might hope that the thinking tax diminishes with more capable models.
We test this hypothesis by comparing Qwen3-8B and Qwen3.5-27B in
thinking mode at identical budgets.
The results are counterintuitive.

\input{../results/paper_figures/table_model_size_scaling}

At $b{=}512$, the 27B model achieves only \textbf{18.3\%} accuracy in
thinking mode---dramatically \emph{worse} than the 8B model's 65.2\%
(Table~\ref{tab:model-size-scaling}).
The thinking tax for 27B ($75.7$\,pp relative to 8B non-thinking at the
same budget) is \textbf{2.6$\times$ larger} than for 8B ($28.8$\,pp).

The root cause is a scaling mismatch between reasoning chain length and
token budget.
The 27B model produces longer, more elaborate reasoning chains---a
consequence of its greater capacity and RLHF training.
At $b{=}512$, the 27B model's natural-stop rate is a mere \textbf{0.7\%}
(Table~\ref{tab:model-size-scaling}, footnote), meaning 99.3\% of its
responses are \emph{truncated projections}: the model ran out of tokens
mid-reasoning and was forced to output an answer.
By contrast, 8B achieves a 37.4\% natural-stop rate at the same budget.

This finding has important practical implications.
As reasoning models grow larger and more capable, their chains of thought
grow proportionally longer, requiring commensurately larger budgets to
avoid truncation.
The thinking tax does not diminish with scale---it \emph{worsens},
making adaptive budget allocation \emph{more} important, not less, for
frontier models.

\begin{figure}[t]
\centering
\includegraphics[width=0.72\textwidth]{fig_thinking_tax_model_size}
\caption{%
  \textbf{The thinking tax scales with model size.}
  Comparison of thinking-mode accuracy for Qwen3-8B (blue) and
  Qwen3.5-27B (red) at matched budgets on GSM8K.
  At $b{=}512$, 27B's accuracy (18.3\%) is 46.9\,pp below 8B's (65.2\%),
  because 27B's longer chains are more severely truncated.
  The 8B non-thinking baseline (green dashed) dominates both at all budgets.
}
\label{fig:thinking-tax-model-size}
\end{figure}

% ---------------------------------------------------------------------------
\subsection{Finding 5: 31.8\% of Questions Are Impossible at All Budgets}
\label{sec:finding:impossible}
% ---------------------------------------------------------------------------

To understand the efficiency ceiling of any budget-allocation scheme, we
categorize GSM8K problems by the \emph{minimum} budget at which Qwen3-8B
solves them correctly in thinking mode.

\input{../results/paper_figures/table_efficiency_frontier}

Table~\ref{tab:efficiency-frontier} and Figure~\ref{fig:difficulty-distribution}
reveal a four-tier difficulty distribution:
\begin{itemize}[nosep,leftmargin=*]
  \item \textbf{Easy} (11.8\%, $n{=}155$): solved at $b{\leq}128$.
    These require only a short chain.
  \item \textbf{Medium} (25.8\%, $n{=}340$): solved at $b{\leq}256$
    but not at 128.  A moderate reasoning chain suffices.
  \item \textbf{Hard} (30.7\%, $n{=}405$): solved only at $b{=}512$.
    These genuinely benefit from extended reasoning.
  \item \textbf{Impossible} (31.8\%, $n{=}419$): unsolved at \emph{any}
    tested budget ($b{\leq}512$).  Tokens spent on these problems
    yield zero return.
\end{itemize}

The ``impossible'' category is surprisingly large: nearly a third of
GSM8K is beyond the model's reach regardless of how much thinking budget
is allocated (up to 512 tokens).
An \emph{oracle router} that (i)~skips impossible questions entirely
and (ii)~assigns the minimum sufficient budget to solvable questions
achieves \textbf{68.2\%} accuracy ($+3.0$\,pp over fixed-512) while
consuming only \textbf{401} average tokens---a \textbf{21.7\%}
savings.\footnote{%
  The oracle accuracy exceeds fixed-512 because it avoids cases where
  extended reasoning causes the model to ``overthink'' and revise a
  correct intermediate answer.}

This analysis establishes both an opportunity and a ceiling for adaptive
methods.
The opportunity: a substantial fraction of compute is wasted on problems
that are either trivially easy (and need few tokens) or genuinely
impossible (and waste all tokens).
The ceiling: no budget controller operating within the $b{\leq}512$
regime can exceed 68.2\% accuracy on this benchmark, because 31.8\% of
problems are fundamentally beyond the model's capability.

\begin{figure}[t]
\centering
\includegraphics[width=0.72\textwidth]{fig3_difficulty_distribution}
\caption{%
  \textbf{Four-tier difficulty distribution of GSM8K problems} under Qwen3-8B thinking mode.
  Problems are categorized by the minimum budget at which they are solved correctly.
  31.8\% of problems are impossible at all tested budgets---tokens spent on these yield zero return.
  An ideal budget allocator would skip impossible questions and assign minimum budgets to the rest.
}
\label{fig:difficulty-distribution}
\end{figure}

\paragraph{Summary of findings.}
Taken together, our five findings paint a coherent picture:
\begin{enumerate}[nosep,leftmargin=*]
  \item The thinking tax is real and large (\S\ref{sec:finding:tax}):
    non-thinking mode dominates at all matched budgets.
  \item Models already signal when thinking was successful
    (\S\ref{sec:finding:natural-stop}): natural stop predicts 96.3\%
    accuracy at zero cost.
  \item High budgets are under-utilized (\S\ref{sec:finding:utilization}):
    up to 60\% of allocated tokens go unused.
  \item The tax worsens with scale (\S\ref{sec:finding:model-size}):
    larger models need proportionally larger budgets.
  \item A third of problems are impossible
    (\S\ref{sec:finding:impossible}): no budget helps.
\end{enumerate}
The natural question is: can we design a system that \emph{thinks only
when thinking actually helps}?
This is precisely the goal of \textsc{Town}, introduced next.

% ===========================================================================

Before introducing our method, we present a systematic empirical study of
the costs and failure modes of chain-of-thought reasoning under constrained
token budgets.  All experiments in this section use
\textbf{Qwen3-8B}~\cite{qwen3_2025} with its native thinking mode
(\texttt{enable\_thinking=True}) and \textbf{Qwen3.5-27B} for cross-scale
validation, evaluated on the full GSM8K test set ($n{=}1{,}319$)
unless otherwise noted.\footnote{%
  A 200-sample subset (seed 42) is used for the budget sweep in
  Table~\ref{tab:thinking-tax}; all other analyses use the complete set.}
We control test-time compute via the \texttt{max\_new\_tokens} parameter,
which caps the \emph{total} number of newly generated tokens---including
both the thinking trace (\texttt{<think>...<\!/think>}) and the final answer.
Crucially, thinking mode and non-thinking mode operate under the
\emph{same} total output-token budget~$b$, ensuring a fair comparison:
\texttt{think@$b$} must fit its entire reasoning chain \emph{plus}
its answer within $b$~tokens, while \texttt{nothink@$b$} uses all
$b$~tokens for the answer directly.
Non-thinking results are obtained by setting
\texttt{enable\_thinking=False} in the Qwen3 chat template, which
suppresses the thinking block entirely.
Greedy decoding ($\tau{=}0$) is used throughout.

Our five findings collectively establish that (i)~thinking mode imposes a
severe \emph{tax} at practical token budgets, (ii)~models already emit
reliable signals indicating when this tax is unnecessary, and
(iii)~a large fraction of compute is wasted on questions the model
cannot solve regardless of budget.
These findings motivate the design of \textsc{Town} in
\S\ref{sec:method}.

% ---------------------------------------------------------------------------
\subsection{Finding 1: Non-Thinking Beats Thinking at All Matched Budgets}
\label{sec:finding:tax}
% ---------------------------------------------------------------------------

Our most striking result is that, when the output-token budget is held
constant, \emph{non-thinking mode uniformly dominates thinking mode}
across every budget level we tested.
Table~\ref{tab:thinking-tax} summarizes the comparison.

% Table 1: The Thinking Tax — Comprehensive comparison at matched budgets
\begin{table}[t]
\centering
\caption{\textbf{The Thinking Tax: Non-thinking mode dramatically outperforms thinking mode at matched token budgets.}
Results on GSM8K with Qwen3-8B.
Rows without a marker use the 200-sample subset (seed=42); starred ($\star$) rows use the full test set ($n{=}1{,}319$).
Qwen3.5-27B results use the full set.
\textbf{At every budget, non-thinking achieves higher accuracy while using fewer tokens.}}
\label{tab:thinking-tax}
\small
\begin{tabular}{ll rrrr}
\toprule
\textbf{Model} & \textbf{Config} & \textbf{Budget} & \textbf{Accuracy} & \textbf{Avg Tok} & \textbf{Early Stop} \\
\midrule
\multirow{10}{*}{\textbf{Qwen3-8B}}
& \texttt{nothink@128}      & 128 & 54.5\%  & 111 & 43.5\% \\
& \texttt{nothink@128$\star$} & 128 & 50.8\%  & 113 & --- \\
& \texttt{nothink@256}      & 256 & 89.0\%  & 140 & 92.0\% \\
& \texttt{nothink@256$\star$} & 256 & \textbf{87.5\%}  & \textbf{146} & 88.8\% \\
& \texttt{nothink@512}      & 512 & 94.0\%  & 145 & 99.5\% \\
\cmidrule{2-6}
& \texttt{think@128}        & 128 & 2.0\%   & 128 & 0.0\% \\
& \texttt{think@128$\star$}  & 128 & 3.0\%   & 128 & --- \\
& \texttt{think@256}        & 256 & 22.0\%  & 255 & 2.0\% \\
& \texttt{think@256$\star$}  & 256 & 18.0\%  & 255 & --- \\
& \texttt{think@512}        & 512 & 66.5\%  & 442 & 47.5\% \\
& \texttt{think@512$\star$}  & 512 & 65.2\%  & 460 & 37.4\% \\
\midrule
\multirow{3}{*}{\textbf{Qwen3.5-27B$\star$}}
& \texttt{think@128}        & 128 & 3.6\%   & 144 & 0.0\% \\
& \texttt{think@256}        & 256 & 7.9\%   & 272 & 0.0\% \\
& \texttt{think@512}        & 512 & 18.3\%  & 528 & 0.7\% \\
\bottomrule
\end{tabular}
\vspace{-2mm}
\end{table}


At budget $b{=}128$, non-thinking achieves \textbf{54.5\%} accuracy while
thinking yields only \textbf{2.0\%}---a 52.5 percentage-point (pp) gap.
At $b{=}256$, the gap narrows but remains enormous:
\textbf{89.0\%} vs.\ \textbf{22.0\%} ($+67.0$\,pp).
Even at $b{=}512$, where thinking has enough room for short chains,
non-thinking still leads by a large margin:
\textbf{94.0\%} vs.\ \textbf{66.5\%} ($+27.5$\,pp).
These results are confirmed on the full GSM8K test set ($n{=}1{,}319$,
$\star$ rows in Table~\ref{tab:thinking-tax}):
\texttt{nothink@256$\star$} achieves \textbf{87.5\%} vs.\ only
\textbf{18.0\%} for \texttt{think@256$\star$}, a
\textbf{69.5\,pp} gap at matched budget.

The root cause is straightforward: a thinking model must first produce a
reasoning chain \emph{inside} its $b$-token budget before emitting an answer.
At low budgets, the chain consumes nearly all available tokens, leaving
insufficient room for a well-formed answer.
A non-thinking model, by contrast, directly generates the solution,
allocating its full budget to answer-bearing content.
This difference manifests in the \emph{Early Stop} column of
Table~\ref{tab:thinking-tax}: at $b{=}256$, 92.0\% of non-thinking
samples terminate naturally before the budget, whereas only 2.0\% of
thinking samples do.
In effect, the non-thinking model has already ``solved'' the problem and
moved on, while the thinking model is still mid-chain.

% Figure 1 (nothink vs thinking) is presented in the Introduction.
% See Figure~\ref{fig:nothink-vs-thinking} for the visual comparison.

We emphasize that this result does \emph{not} imply that chain-of-thought
reasoning is useless---at sufficiently large budgets (e.g.,\ $b{=}4096$),
thinking mode recovers and eventually surpasses non-thinking.
Rather, it shows that under the budget-constrained regime that is
increasingly relevant for cost-sensitive deployment, the overhead of
producing a reasoning chain is a \emph{tax} that must be justified by
the problem's difficulty.

% ---------------------------------------------------------------------------
\subsection{Finding 2: Natural Stop Is a Free Confidence Oracle}
\label{sec:finding:natural-stop}
% ---------------------------------------------------------------------------

When a thinking model generates a complete reasoning chain and emits
a ``Final answer'' marker \emph{before} exhausting its token budget,
we say the sample exhibits a \emph{natural stop}.
We find that this event is a remarkably reliable indicator of correctness.

% Table 1: Natural Stop as Confidence Oracle
% Key result: natural-stop samples achieve 99.0% accuracy at budget=512 (HF engine)
% Updated 2026-04-09: switched from vLLM+projection to HF engine for consistency
\begin{table}[t]
\centering
\caption{\textbf{Natural stop as a free confidence oracle.}
When a thinking model naturally terminates before exhausting its token budget,
the resulting answer is correct with remarkably high probability.
This phenomenon is consistent across models and benchmarks.
\emph{NatStop}: fraction of samples that stop early ($<$95\% of budget used).
\emph{Acc$_\text{nat}$}: accuracy among early-stop samples.
\emph{Acc$_\text{hit}$}: accuracy among budget-exhausting samples.
All Qwen3-8B results use HF engine with greedy decoding.
DeepSeek-R1 Acc splits are omitted due to small sample sizes ($n{=}40$--$80$).}
\label{tab:natural-stop-oracle}
\small
\setlength{\tabcolsep}{3pt}
\begin{tabular}{@{}lrrrrrr@{}}
\toprule
\textbf{Model / Benchmark} & \textbf{Bud.} & \textbf{Acc} & \textbf{NatStop} & \textbf{Acc$_\text{nat}$} & \textbf{Acc$_\text{hit}$} & \textbf{Gap} \\
\midrule
\multirow{3}{*}{Qwen3-8B / GSM8K (n=1319)}
 & 128  &  3.0 &  0.0 & --- &  3.0 & --- \\
 & 256  & 18.0 &  1.4 & 100.0 & 16.8 & +83.2 \\
 & 512  & 56.9 & 37.4 & 99.0 & 31.8 & \textbf{+67.2} \\
\midrule
\multirow{3}{*}{DeepSeek-R1-8B / GSM8K (n=40)}
 & 256  & 22.5 & 15.0 & --- & --- & --- \\
 & 512  & 55.0 & 85.0 & --- & --- & --- \\
 & 1024 & 65.0 & 97.5 & --- & --- & --- \\
\midrule
\multirow{3}{*}{DeepSeek-R1-8B / MATH500 (n=80)}
 & 1024 & 25.0 & 31.2 & --- & --- & --- \\
 & 2048 & 38.7 & 55.0 & --- & --- & --- \\
 & 4096 & 46.2 & 76.2 & --- & --- & --- \\
\bottomrule
\end{tabular}
\end{table}


At $b{=}512$ with Qwen3-8B, 37.4\% of samples terminate before consuming
95\% of their token budget---we call these \emph{natural stops}
(Table~\ref{tab:natural-stop-oracle}).
Among these 475 samples, accuracy is \textbf{96.3\%}---nearly perfect.
Among samples that exhaust the budget (i.e., are truncated), accuracy is
only 47.5\%, yielding an accuracy gap of \textbf{$+$48.8\,pp}
(Figure~\ref{fig:natural-stop-split}).
A stricter definition---requiring an explicit ``Final answer'' marker
(\texttt{has\_final})---selects only 6.1\% of samples but achieves
93.8\% accuracy, further confirming the signal's reliability.

This phenomenon admits a simple explanation.
A natural stop means the model's \emph{own} generation process
determined that the reasoning chain was complete---the model ``believes''
it has reached a satisfactory answer.
While LLM confidence is notoriously miscalibrated in
general~\cite{chen2024overthinking}, the natural-stop signal is
structurally different from soft confidence scores: it is a
\emph{binary}, \emph{endogenous} event that does not require access to
logits, hidden states, or any post-hoc calibration.
It is, in a precise sense, a \emph{free} confidence oracle---one that
the model provides at zero additional inference cost.

\begin{remark}[Generality]
The natural-stop oracle effect is not limited to Qwen3-8B.
Table~\ref{tab:natural-stop-oracle} shows that DeepSeek-R1-8B exhibits
similar behavior: at $b{=}512$, 85.0\% of GSM8K samples stop naturally,
and on the harder MATH500 benchmark, natural-stop rates increase steadily
with budget (31.2\% $\to$ 55.0\% $\to$ 76.2\% as $b$ scales from 1024
to 4096).
This universality suggests that the phenomenon is an intrinsic property
of reasoning-trained models, not an artifact of a specific architecture.
\end{remark}

\begin{figure}[t]
\centering
\includegraphics[width=0.72\textwidth]{fig9_acc_split_natural_vs_hit}
\caption{%
  \textbf{Natural stop is a free confidence oracle.}
  Accuracy decomposition by stopping behavior for Qwen3-8B on GSM8K.
  Samples that stop naturally before the budget (blue) achieve dramatically
  higher accuracy than those that exhaust the budget (red), across all
  budget levels.  At $b{=}512$, natural-stop samples reach 96.3\% accuracy
  vs.\ 47.5\% for truncated samples---a 48.8\,pp gap.
}
\label{fig:natural-stop-split}
\end{figure}

% ---------------------------------------------------------------------------
\subsection{Finding 3: Token Utilization Collapses at High Budgets}
\label{sec:finding:utilization}
% ---------------------------------------------------------------------------

A complementary perspective on the natural-stop phenomenon comes from
examining \emph{token utilization}: the fraction of the allocated budget
that the model actually consumes.

% Table 2: Token Utilization Decreases at Higher Budgets
% Key result: utilization drops from 100% to 38% as budget increases
\begin{table}[t]
\centering
\caption{\textbf{Token utilization ratio decreases with budget.}
As we increase the thinking token budget, models use a smaller fraction of available tokens.
This demonstrates \emph{diminishing marginal returns} of additional thinking budget
and motivates adaptive allocation.
Utilization = avg tokens used / budget.
$^{\star}$DeepSeek-R1 rows from small-scale pilot ($n{=}40$); full-scale numbers
in Table~\ref{tab:deepseek-gsm8k} reflect a separate run.}
\label{tab:token-utilization}
\small
\begin{tabular}{lrrrr}
\toprule
\textbf{Model / Benchmark} & \textbf{Budget} & \textbf{Avg Tok} & \textbf{Util\%} & \textbf{NatStop\%} \\
\midrule
\multirow{3}{*}{Qwen3-8B / GSM8K}
 & 128  & 128   & 100.0 &  0.0 \\
 & 256  & 255   &  99.7 &  1.4 \\
 & 512  & 460   &  89.8 & 37.4 \\
\midrule
\multirow{3}{*}{DeepSeek-R1 / GSM8K$^{\star}$}
 & 256  & 264   & 103.0 & 15.0 \\
 & 512  & 376   &  73.5 & 85.0 \\
 & 1024 & 393   &  \textbf{38.3} & 97.5 \\
\midrule
\multirow{3}{*}{DeepSeek-R1 / MATH500}
 & 1024 & 910   &  88.8 & 31.2 \\
 & 2048 & 1551  &  75.7 & 55.0 \\
 & 4096 & 2353  &  57.4 & 76.2 \\
\bottomrule
\end{tabular}
\vspace{-2mm}
\end{table}


Table~\ref{tab:token-utilization} reveals a striking pattern.
At $b{=}128$, Qwen3-8B uses 100\% of its budget---every token is
consumed, and the model is almost certainly still mid-reasoning when
generation is forcibly terminated.
At $b{=}256$, utilization remains near-perfect (99.7\%), with only 1.4\%
of samples stopping naturally.
But at $b{=}512$, utilization drops to \textbf{89.8\%}, with an average
of only 460 tokens consumed out of 512 available.
The model is already ``paying'' for $\sim$52 tokens it does not need.

The pattern is even more dramatic for DeepSeek-R1 on GSM8K: utilization
plummets from 103.0\% at $b{=}256$ (the model slightly exceeds its
budget due to token-counting granularity) to \textbf{38.3\%} at
$b{=}1{,}024$, with 97.5\% of samples stopping naturally.
This means the model uses fewer than 400 tokens on average when allocated
1{,}024---more than 60\% of the budget is wasted.

This utilization collapse is a direct consequence of the heterogeneous
difficulty distribution in typical benchmarks: many problems are easy
enough that a short reasoning chain suffices, and additional budget
yields no benefit.
It provides strong empirical motivation for adaptive budget allocation:
rather than setting a single high budget and accepting low utilization,
one should match the budget to the problem's actual difficulty.

% ---------------------------------------------------------------------------
\subsection{Finding 4: The Thinking Tax Scales with Model Size}
\label{sec:finding:model-size}
% ---------------------------------------------------------------------------

One might hope that the thinking tax diminishes with more capable models.
We test this hypothesis by comparing Qwen3-8B and Qwen3.5-27B in
thinking mode at identical budgets.
The results are counterintuitive.

% Table: Thinking Tax scales with model size (8B, 9B, 27B)
% Updated 2026-04-08 with gap-fill results: 8B nothink@512/1024/2048, 9B think@512/nothink@256/512, 8B think@1024
\begin{table}[t]
\centering
\caption{\textbf{The Thinking Tax worsens with model size} on GSM8K ($n{=}1{,}319$).
At budget=512, thinking-mode accuracy collapses from 56.9\% (8B) to 18.3\% (27B),
while non-thinking accuracy remains high across all scales:
8B=93.1\%, 9B=93.2\%, 27B=\textbf{95.5\%}.
The 9B model exhibits the largest tax (\textbf{77.7\,pp}), confirming that
inverse scaling is driven by chain-length growth, not capability loss.}
\label{tab:model-size-scaling}
\small
\begin{tabular}{l rr rr rr}
\toprule
 & \multicolumn{2}{c}{\textbf{8B}} & \multicolumn{2}{c}{\textbf{9B}} & \multicolumn{2}{c}{\textbf{27B}} \\
\cmidrule(lr){2-3} \cmidrule(lr){4-5} \cmidrule(lr){6-7}
\textbf{Budget} & \textbf{Think} & \textbf{NoThink} & \textbf{Think} & \textbf{NoThink} & \textbf{Think} & \textbf{NoThink} \\
\midrule
128 & 3.0\%  & 50.8\% & ---    & ---    & 3.6\%  & 9.9\% \\
256 & 18.0\% & 87.5\% & 3.4\%  & 61.2\% & 7.9\%  & 65.1\% \\
512 & 56.9\% & 93.1\% & 15.5\% & 93.2\% & 18.3\% & \textbf{95.5\%} \\
1024 & 86.1\% & 93.1\% & 41.8\% & ---$^\ddagger$ & ${\sim}$49\%$^\dagger$ & --- \\
2048 & ---$^\ddagger$ & 93.1\% & 66.8\% & ---    & ---    & --- \\
\midrule
\multicolumn{7}{l}{\textit{Thinking Tax at budget=512 (NoThink@512 $-$ Think@512):}} \\
\midrule
Tax & \multicolumn{6}{c}{8B: \textbf{36.2\,pp} (93.1\%$-$56.9\%); \quad 9B: \textbf{77.7\,pp} (93.2\%$-$15.5\%); \quad 27B: \textbf{77.2\,pp} (95.5\%$-$18.3\%)} \\
\bottomrule
\end{tabular}
\vspace{1mm}

\footnotesize
$^\dagger$Checkpoint at $n{=}200$; full-dataset experiment pending.\\
$^\dagger$Checkpoint at $n{=}200$; full-dataset experiment pending.
$^\ddagger$8B think@2048 in progress.\\
\textit{Root cause}: natural-stop rates at budget=512: 8B=37.4\%, 9B${\approx}$0.5\%, 27B=0.7\%.\\
At $b{=}1024$: 8B think=86.1\% vs nothink=93.1\% (gap 7.0\,pp); crossover not yet reached.\\
Nothink saturates: 8B nothink@512/1024/2048 = 93.1\%/93.1\%/93.1\% (avg ${\sim}$152 tokens).\\
\vspace{-3mm}
\end{table}


At $b{=}512$, the 27B model achieves only \textbf{18.3\%} accuracy in
thinking mode---dramatically \emph{worse} than the 8B model's 65.2\%
(Table~\ref{tab:model-size-scaling}).
The thinking tax for 27B ($75.7$\,pp relative to 8B non-thinking at the
same budget) is \textbf{2.6$\times$ larger} than for 8B ($28.8$\,pp).

The root cause is a scaling mismatch between reasoning chain length and
token budget.
The 27B model produces longer, more elaborate reasoning chains---a
consequence of its greater capacity and RLHF training.
At $b{=}512$, the 27B model's natural-stop rate is a mere \textbf{0.7\%}
(Table~\ref{tab:model-size-scaling}, footnote), meaning 99.3\% of its
responses are \emph{truncated projections}: the model ran out of tokens
mid-reasoning and was forced to output an answer.
By contrast, 8B achieves a 37.4\% natural-stop rate at the same budget.

This finding has important practical implications.
As reasoning models grow larger and more capable, their chains of thought
grow proportionally longer, requiring commensurately larger budgets to
avoid truncation.
The thinking tax does not diminish with scale---it \emph{worsens},
making adaptive budget allocation \emph{more} important, not less, for
frontier models.

\begin{figure}[t]
\centering
\includegraphics[width=0.72\textwidth]{fig_thinking_tax_model_size}
\caption{%
  \textbf{The thinking tax scales with model size.}
  Comparison of thinking-mode accuracy for Qwen3-8B (blue) and
  Qwen3.5-27B (red) at matched budgets on GSM8K.
  At $b{=}512$, 27B's accuracy (18.3\%) is 46.9\,pp below 8B's (65.2\%),
  because 27B's longer chains are more severely truncated.
  The 8B non-thinking baseline (green dashed) dominates both at all budgets.
}
\label{fig:thinking-tax-model-size}
\end{figure}

% ---------------------------------------------------------------------------
\subsection{Finding 5: 31.8\% of Questions Are Impossible at All Budgets}
\label{sec:finding:impossible}
% ---------------------------------------------------------------------------

To understand the efficiency ceiling of any budget-allocation scheme, we
categorize GSM8K problems by the \emph{minimum} budget at which Qwen3-8B
solves them correctly in thinking mode.

% Table 3: Thinking Efficiency Frontier
% Key result: 31.8% of questions are "impossible" at any tested budget
\begin{table}[t]
\centering
\caption{\textbf{Thinking Efficiency Frontier: difficulty distribution.}
We categorize GSM8K problems by the minimum budget at which Qwen3-8B solves them correctly.
31.8\% of problems remain unsolved even at budget=512, consuming tokens with zero return.
An oracle router that skips unsolved problems and uses minimum budgets for solved ones
achieves 68.2\% accuracy ($+$3.0pp over fixed-512) while using only 401 avg tokens (21.7\% savings).}
\label{tab:efficiency-frontier}
\small
\begin{tabular}{lrrrrr}
\toprule
\textbf{Category} & \textbf{Min Budget} & \textbf{Count} & \textbf{\%} & \textbf{Cumul. Acc} & \textbf{$\Delta$ Acc} \\
\midrule
Easy       & $\leq$128 & 155  & 11.8 & 11.8 & +11.8 \\
Medium     & $\leq$256 & 340  & 25.8 & 37.5 & +25.8 \\
Hard       & $\leq$512 & 405  & 30.7 & 68.2 & +30.7 \\
Impossible & ---       & 419  & 31.8 & ---  & --- \\
\midrule
\multicolumn{2}{l}{\textbf{Total}} & \textbf{1319} & 100 & & \\
\bottomrule
\end{tabular}
\end{table}


Table~\ref{tab:efficiency-frontier} and Figure~\ref{fig:difficulty-distribution}
reveal a four-tier difficulty distribution:
\begin{itemize}[nosep,leftmargin=*]
  \item \textbf{Easy} (11.8\%, $n{=}155$): solved at $b{\leq}128$.
    These require only a short chain.
  \item \textbf{Medium} (25.8\%, $n{=}340$): solved at $b{\leq}256$
    but not at 128.  A moderate reasoning chain suffices.
  \item \textbf{Hard} (30.7\%, $n{=}405$): solved only at $b{=}512$.
    These genuinely benefit from extended reasoning.
  \item \textbf{Impossible} (31.8\%, $n{=}419$): unsolved at \emph{any}
    tested budget ($b{\leq}512$).  Tokens spent on these problems
    yield zero return.
\end{itemize}

The ``impossible'' category is surprisingly large: nearly a third of
GSM8K is beyond the model's reach regardless of how much thinking budget
is allocated (up to 512 tokens).
An \emph{oracle router} that (i)~skips impossible questions entirely
and (ii)~assigns the minimum sufficient budget to solvable questions
achieves \textbf{68.2\%} accuracy ($+3.0$\,pp over fixed-512) while
consuming only \textbf{401} average tokens---a \textbf{21.7\%}
savings.\footnote{%
  The oracle accuracy exceeds fixed-512 because it avoids cases where
  extended reasoning causes the model to ``overthink'' and revise a
  correct intermediate answer.}

This analysis establishes both an opportunity and a ceiling for adaptive
methods.
The opportunity: a substantial fraction of compute is wasted on problems
that are either trivially easy (and need few tokens) or genuinely
impossible (and waste all tokens).
The ceiling: no budget controller operating within the $b{\leq}512$
regime can exceed 68.2\% accuracy on this benchmark, because 31.8\% of
problems are fundamentally beyond the model's capability.

\begin{figure}[t]
\centering
\includegraphics[width=0.72\textwidth]{fig3_difficulty_distribution}
\caption{%
  \textbf{Four-tier difficulty distribution of GSM8K problems} under Qwen3-8B thinking mode.
  Problems are categorized by the minimum budget at which they are solved correctly.
  31.8\% of problems are impossible at all tested budgets---tokens spent on these yield zero return.
  An ideal budget allocator would skip impossible questions and assign minimum budgets to the rest.
}
\label{fig:difficulty-distribution}
\end{figure}

\paragraph{Summary of findings.}
Taken together, our five findings paint a coherent picture:
\begin{enumerate}[nosep,leftmargin=*]
  \item The thinking tax is real and large (\S\ref{sec:finding:tax}):
    non-thinking mode dominates at all matched budgets.
  \item Models already signal when thinking was successful
    (\S\ref{sec:finding:natural-stop}): natural stop predicts 96.3\%
    accuracy at zero cost.
  \item High budgets are under-utilized (\S\ref{sec:finding:utilization}):
    up to 60\% of allocated tokens go unused.
  \item The tax worsens with scale (\S\ref{sec:finding:model-size}):
    larger models need proportionally larger budgets.
  \item A third of problems are impossible
    (\S\ref{sec:finding:impossible}): no budget helps.
\end{enumerate}
The natural question is: can we design a system that \emph{thinks only
when thinking actually helps}?
This is precisely the goal of \textsc{Town}, introduced next.

% ===========================================================================

Before introducing our method, we present a systematic empirical study of
the costs and failure modes of chain-of-thought reasoning under constrained
token budgets.  All experiments in this section use
\textbf{Qwen3-8B}~\cite{qwen3_2025} with its native thinking mode
(\texttt{enable\_thinking=True}) and \textbf{Qwen3.5-27B} for cross-scale
validation, evaluated on the full GSM8K test set ($n{=}1{,}319$)
unless otherwise noted.\footnote{%
  A 200-sample subset (seed 42) is used for the budget sweep in
  Table~\ref{tab:thinking-tax}; all other analyses use the complete set.}
We control test-time compute via the \texttt{max\_new\_tokens} parameter,
which caps the \emph{total} number of newly generated tokens---including
both the thinking trace (\texttt{<think>...<\!/think>}) and the final answer.
Crucially, thinking mode and non-thinking mode operate under the
\emph{same} total output-token budget~$b$, ensuring a fair comparison:
\texttt{think@$b$} must fit its entire reasoning chain \emph{plus}
its answer within $b$~tokens, while \texttt{nothink@$b$} uses all
$b$~tokens for the answer directly.
Non-thinking results are obtained by setting
\texttt{enable\_thinking=False} in the Qwen3 chat template, which
suppresses the thinking block entirely.
Greedy decoding ($\tau{=}0$) is used throughout.

Our five findings collectively establish that (i)~thinking mode imposes a
severe \emph{tax} at practical token budgets, (ii)~models already emit
reliable signals indicating when this tax is unnecessary, and
(iii)~a large fraction of compute is wasted on questions the model
cannot solve regardless of budget.
These findings motivate the design of \textsc{Town} in
\S\ref{sec:method}.

% ---------------------------------------------------------------------------
\subsection{Finding 1: Non-Thinking Beats Thinking at All Matched Budgets}
\label{sec:finding:tax}
% ---------------------------------------------------------------------------

Our most striking result is that, when the output-token budget is held
constant, \emph{non-thinking mode uniformly dominates thinking mode}
across every budget level we tested.
Table~\ref{tab:thinking-tax} summarizes the comparison.

% Table 1: The Thinking Tax — Comprehensive comparison at matched budgets
\begin{table}[t]
\centering
\caption{\textbf{The Thinking Tax: Non-thinking mode dramatically outperforms thinking mode at matched token budgets.}
Results on GSM8K with Qwen3-8B.
Rows without a marker use the 200-sample subset (seed=42); starred ($\star$) rows use the full test set ($n{=}1{,}319$).
Qwen3.5-27B results use the full set.
\textbf{At every budget, non-thinking achieves higher accuracy while using fewer tokens.}}
\label{tab:thinking-tax}
\small
\begin{tabular}{ll rrrr}
\toprule
\textbf{Model} & \textbf{Config} & \textbf{Budget} & \textbf{Accuracy} & \textbf{Avg Tok} & \textbf{Early Stop} \\
\midrule
\multirow{10}{*}{\textbf{Qwen3-8B}}
& \texttt{nothink@128}      & 128 & 54.5\%  & 111 & 43.5\% \\
& \texttt{nothink@128$\star$} & 128 & 50.8\%  & 113 & --- \\
& \texttt{nothink@256}      & 256 & 89.0\%  & 140 & 92.0\% \\
& \texttt{nothink@256$\star$} & 256 & \textbf{87.5\%}  & \textbf{146} & 88.8\% \\
& \texttt{nothink@512}      & 512 & 94.0\%  & 145 & 99.5\% \\
\cmidrule{2-6}
& \texttt{think@128}        & 128 & 2.0\%   & 128 & 0.0\% \\
& \texttt{think@128$\star$}  & 128 & 3.0\%   & 128 & --- \\
& \texttt{think@256}        & 256 & 22.0\%  & 255 & 2.0\% \\
& \texttt{think@256$\star$}  & 256 & 18.0\%  & 255 & --- \\
& \texttt{think@512}        & 512 & 66.5\%  & 442 & 47.5\% \\
& \texttt{think@512$\star$}  & 512 & 65.2\%  & 460 & 37.4\% \\
\midrule
\multirow{3}{*}{\textbf{Qwen3.5-27B$\star$}}
& \texttt{think@128}        & 128 & 3.6\%   & 144 & 0.0\% \\
& \texttt{think@256}        & 256 & 7.9\%   & 272 & 0.0\% \\
& \texttt{think@512}        & 512 & 18.3\%  & 528 & 0.7\% \\
\bottomrule
\end{tabular}
\vspace{-2mm}
\end{table}


At budget $b{=}128$, non-thinking achieves \textbf{54.5\%} accuracy while
thinking yields only \textbf{2.0\%}---a 52.5 percentage-point (pp) gap.
At $b{=}256$, the gap narrows but remains enormous:
\textbf{89.0\%} vs.\ \textbf{22.0\%} ($+67.0$\,pp).
Even at $b{=}512$, where thinking has enough room for short chains,
non-thinking still leads by a large margin:
\textbf{94.0\%} vs.\ \textbf{66.5\%} ($+27.5$\,pp).
These results are confirmed on the full GSM8K test set ($n{=}1{,}319$,
$\star$ rows in Table~\ref{tab:thinking-tax}):
\texttt{nothink@256$\star$} achieves \textbf{87.5\%} vs.\ only
\textbf{18.0\%} for \texttt{think@256$\star$}, a
\textbf{69.5\,pp} gap at matched budget.

The root cause is straightforward: a thinking model must first produce a
reasoning chain \emph{inside} its $b$-token budget before emitting an answer.
At low budgets, the chain consumes nearly all available tokens, leaving
insufficient room for a well-formed answer.
A non-thinking model, by contrast, directly generates the solution,
allocating its full budget to answer-bearing content.
This difference manifests in the \emph{Early Stop} column of
Table~\ref{tab:thinking-tax}: at $b{=}256$, 92.0\% of non-thinking
samples terminate naturally before the budget, whereas only 2.0\% of
thinking samples do.
In effect, the non-thinking model has already ``solved'' the problem and
moved on, while the thinking model is still mid-chain.

% Figure 1 (nothink vs thinking) is presented in the Introduction.
% See Figure~\ref{fig:nothink-vs-thinking} for the visual comparison.

We emphasize that this result does \emph{not} imply that chain-of-thought
reasoning is useless---at sufficiently large budgets (e.g.,\ $b{=}4096$),
thinking mode recovers and eventually surpasses non-thinking.
Rather, it shows that under the budget-constrained regime that is
increasingly relevant for cost-sensitive deployment, the overhead of
producing a reasoning chain is a \emph{tax} that must be justified by
the problem's difficulty.

% ---------------------------------------------------------------------------
\subsection{Finding 2: Natural Stop Is a Free Confidence Oracle}
\label{sec:finding:natural-stop}
% ---------------------------------------------------------------------------

When a thinking model generates a complete reasoning chain and emits
a ``Final answer'' marker \emph{before} exhausting its token budget,
we say the sample exhibits a \emph{natural stop}.
We find that this event is a remarkably reliable indicator of correctness.

% Table 1: Natural Stop as Confidence Oracle
% Key result: natural-stop samples achieve 99.0% accuracy at budget=512 (HF engine)
% Updated 2026-04-09: switched from vLLM+projection to HF engine for consistency
\begin{table}[t]
\centering
\caption{\textbf{Natural stop as a free confidence oracle.}
When a thinking model naturally terminates before exhausting its token budget,
the resulting answer is correct with remarkably high probability.
This phenomenon is consistent across models and benchmarks.
\emph{NatStop}: fraction of samples that stop early ($<$95\% of budget used).
\emph{Acc$_\text{nat}$}: accuracy among early-stop samples.
\emph{Acc$_\text{hit}$}: accuracy among budget-exhausting samples.
All Qwen3-8B results use HF engine with greedy decoding.
DeepSeek-R1 Acc splits are omitted due to small sample sizes ($n{=}40$--$80$).}
\label{tab:natural-stop-oracle}
\small
\setlength{\tabcolsep}{3pt}
\begin{tabular}{@{}lrrrrrr@{}}
\toprule
\textbf{Model / Benchmark} & \textbf{Bud.} & \textbf{Acc} & \textbf{NatStop} & \textbf{Acc$_\text{nat}$} & \textbf{Acc$_\text{hit}$} & \textbf{Gap} \\
\midrule
\multirow{3}{*}{Qwen3-8B / GSM8K (n=1319)}
 & 128  &  3.0 &  0.0 & --- &  3.0 & --- \\
 & 256  & 18.0 &  1.4 & 100.0 & 16.8 & +83.2 \\
 & 512  & 56.9 & 37.4 & 99.0 & 31.8 & \textbf{+67.2} \\
\midrule
\multirow{3}{*}{DeepSeek-R1-8B / GSM8K (n=40)}
 & 256  & 22.5 & 15.0 & --- & --- & --- \\
 & 512  & 55.0 & 85.0 & --- & --- & --- \\
 & 1024 & 65.0 & 97.5 & --- & --- & --- \\
\midrule
\multirow{3}{*}{DeepSeek-R1-8B / MATH500 (n=80)}
 & 1024 & 25.0 & 31.2 & --- & --- & --- \\
 & 2048 & 38.7 & 55.0 & --- & --- & --- \\
 & 4096 & 46.2 & 76.2 & --- & --- & --- \\
\bottomrule
\end{tabular}
\end{table}


At $b{=}512$ with Qwen3-8B, 37.4\% of samples terminate before consuming
95\% of their token budget---we call these \emph{natural stops}
(Table~\ref{tab:natural-stop-oracle}).
Among these 475 samples, accuracy is \textbf{96.3\%}---nearly perfect.
Among samples that exhaust the budget (i.e., are truncated), accuracy is
only 47.5\%, yielding an accuracy gap of \textbf{$+$48.8\,pp}
(Figure~\ref{fig:natural-stop-split}).
A stricter definition---requiring an explicit ``Final answer'' marker
(\texttt{has\_final})---selects only 6.1\% of samples but achieves
93.8\% accuracy, further confirming the signal's reliability.

This phenomenon admits a simple explanation.
A natural stop means the model's \emph{own} generation process
determined that the reasoning chain was complete---the model ``believes''
it has reached a satisfactory answer.
While LLM confidence is notoriously miscalibrated in
general~\cite{chen2024overthinking}, the natural-stop signal is
structurally different from soft confidence scores: it is a
\emph{binary}, \emph{endogenous} event that does not require access to
logits, hidden states, or any post-hoc calibration.
It is, in a precise sense, a \emph{free} confidence oracle---one that
the model provides at zero additional inference cost.

\begin{remark}[Generality]
The natural-stop oracle effect is not limited to Qwen3-8B.
Table~\ref{tab:natural-stop-oracle} shows that DeepSeek-R1-8B exhibits
similar behavior: at $b{=}512$, 85.0\% of GSM8K samples stop naturally,
and on the harder MATH500 benchmark, natural-stop rates increase steadily
with budget (31.2\% $\to$ 55.0\% $\to$ 76.2\% as $b$ scales from 1024
to 4096).
This universality suggests that the phenomenon is an intrinsic property
of reasoning-trained models, not an artifact of a specific architecture.
\end{remark}

\begin{figure}[t]
\centering
\includegraphics[width=0.72\textwidth]{fig9_acc_split_natural_vs_hit}
\caption{%
  \textbf{Natural stop is a free confidence oracle.}
  Accuracy decomposition by stopping behavior for Qwen3-8B on GSM8K.
  Samples that stop naturally before the budget (blue) achieve dramatically
  higher accuracy than those that exhaust the budget (red), across all
  budget levels.  At $b{=}512$, natural-stop samples reach 96.3\% accuracy
  vs.\ 47.5\% for truncated samples---a 48.8\,pp gap.
}
\label{fig:natural-stop-split}
\end{figure}

% ---------------------------------------------------------------------------
\subsection{Finding 3: Token Utilization Collapses at High Budgets}
\label{sec:finding:utilization}
% ---------------------------------------------------------------------------

A complementary perspective on the natural-stop phenomenon comes from
examining \emph{token utilization}: the fraction of the allocated budget
that the model actually consumes.

% Table 2: Token Utilization Decreases at Higher Budgets
% Key result: utilization drops from 100% to 38% as budget increases
\begin{table}[t]
\centering
\caption{\textbf{Token utilization ratio decreases with budget.}
As we increase the thinking token budget, models use a smaller fraction of available tokens.
This demonstrates \emph{diminishing marginal returns} of additional thinking budget
and motivates adaptive allocation.
Utilization = avg tokens used / budget.
$^{\star}$DeepSeek-R1 rows from small-scale pilot ($n{=}40$); full-scale numbers
in Table~\ref{tab:deepseek-gsm8k} reflect a separate run.}
\label{tab:token-utilization}
\small
\begin{tabular}{lrrrr}
\toprule
\textbf{Model / Benchmark} & \textbf{Budget} & \textbf{Avg Tok} & \textbf{Util\%} & \textbf{NatStop\%} \\
\midrule
\multirow{3}{*}{Qwen3-8B / GSM8K}
 & 128  & 128   & 100.0 &  0.0 \\
 & 256  & 255   &  99.7 &  1.4 \\
 & 512  & 460   &  89.8 & 37.4 \\
\midrule
\multirow{3}{*}{DeepSeek-R1 / GSM8K$^{\star}$}
 & 256  & 264   & 103.0 & 15.0 \\
 & 512  & 376   &  73.5 & 85.0 \\
 & 1024 & 393   &  \textbf{38.3} & 97.5 \\
\midrule
\multirow{3}{*}{DeepSeek-R1 / MATH500}
 & 1024 & 910   &  88.8 & 31.2 \\
 & 2048 & 1551  &  75.7 & 55.0 \\
 & 4096 & 2353  &  57.4 & 76.2 \\
\bottomrule
\end{tabular}
\vspace{-2mm}
\end{table}


Table~\ref{tab:token-utilization} reveals a striking pattern.
At $b{=}128$, Qwen3-8B uses 100\% of its budget---every token is
consumed, and the model is almost certainly still mid-reasoning when
generation is forcibly terminated.
At $b{=}256$, utilization remains near-perfect (99.7\%), with only 1.4\%
of samples stopping naturally.
But at $b{=}512$, utilization drops to \textbf{89.8\%}, with an average
of only 460 tokens consumed out of 512 available.
The model is already ``paying'' for $\sim$52 tokens it does not need.

The pattern is even more dramatic for DeepSeek-R1 on GSM8K: utilization
plummets from 103.0\% at $b{=}256$ (the model slightly exceeds its
budget due to token-counting granularity) to \textbf{38.3\%} at
$b{=}1{,}024$, with 97.5\% of samples stopping naturally.
This means the model uses fewer than 400 tokens on average when allocated
1{,}024---more than 60\% of the budget is wasted.

This utilization collapse is a direct consequence of the heterogeneous
difficulty distribution in typical benchmarks: many problems are easy
enough that a short reasoning chain suffices, and additional budget
yields no benefit.
It provides strong empirical motivation for adaptive budget allocation:
rather than setting a single high budget and accepting low utilization,
one should match the budget to the problem's actual difficulty.

% ---------------------------------------------------------------------------
\subsection{Finding 4: The Thinking Tax Scales with Model Size}
\label{sec:finding:model-size}
% ---------------------------------------------------------------------------

One might hope that the thinking tax diminishes with more capable models.
We test this hypothesis by comparing Qwen3-8B and Qwen3.5-27B in
thinking mode at identical budgets.
The results are counterintuitive.

% Table: Thinking Tax scales with model size (8B, 9B, 27B)
% Updated 2026-04-08 with gap-fill results: 8B nothink@512/1024/2048, 9B think@512/nothink@256/512, 8B think@1024
\begin{table}[t]
\centering
\caption{\textbf{The Thinking Tax worsens with model size} on GSM8K ($n{=}1{,}319$).
At budget=512, thinking-mode accuracy collapses from 56.9\% (8B) to 18.3\% (27B),
while non-thinking accuracy remains high across all scales:
8B=93.1\%, 9B=93.2\%, 27B=\textbf{95.5\%}.
The 9B model exhibits the largest tax (\textbf{77.7\,pp}), confirming that
inverse scaling is driven by chain-length growth, not capability loss.}
\label{tab:model-size-scaling}
\small
\begin{tabular}{l rr rr rr}
\toprule
 & \multicolumn{2}{c}{\textbf{8B}} & \multicolumn{2}{c}{\textbf{9B}} & \multicolumn{2}{c}{\textbf{27B}} \\
\cmidrule(lr){2-3} \cmidrule(lr){4-5} \cmidrule(lr){6-7}
\textbf{Budget} & \textbf{Think} & \textbf{NoThink} & \textbf{Think} & \textbf{NoThink} & \textbf{Think} & \textbf{NoThink} \\
\midrule
128 & 3.0\%  & 50.8\% & ---    & ---    & 3.6\%  & 9.9\% \\
256 & 18.0\% & 87.5\% & 3.4\%  & 61.2\% & 7.9\%  & 65.1\% \\
512 & 56.9\% & 93.1\% & 15.5\% & 93.2\% & 18.3\% & \textbf{95.5\%} \\
1024 & 86.1\% & 93.1\% & 41.8\% & ---$^\ddagger$ & ${\sim}$49\%$^\dagger$ & --- \\
2048 & ---$^\ddagger$ & 93.1\% & 66.8\% & ---    & ---    & --- \\
\midrule
\multicolumn{7}{l}{\textit{Thinking Tax at budget=512 (NoThink@512 $-$ Think@512):}} \\
\midrule
Tax & \multicolumn{6}{c}{8B: \textbf{36.2\,pp} (93.1\%$-$56.9\%); \quad 9B: \textbf{77.7\,pp} (93.2\%$-$15.5\%); \quad 27B: \textbf{77.2\,pp} (95.5\%$-$18.3\%)} \\
\bottomrule
\end{tabular}
\vspace{1mm}

\footnotesize
$^\dagger$Checkpoint at $n{=}200$; full-dataset experiment pending.\\
$^\dagger$Checkpoint at $n{=}200$; full-dataset experiment pending.
$^\ddagger$8B think@2048 in progress.\\
\textit{Root cause}: natural-stop rates at budget=512: 8B=37.4\%, 9B${\approx}$0.5\%, 27B=0.7\%.\\
At $b{=}1024$: 8B think=86.1\% vs nothink=93.1\% (gap 7.0\,pp); crossover not yet reached.\\
Nothink saturates: 8B nothink@512/1024/2048 = 93.1\%/93.1\%/93.1\% (avg ${\sim}$152 tokens).\\
\vspace{-3mm}
\end{table}


At $b{=}512$, the 27B model achieves only \textbf{18.3\%} accuracy in
thinking mode---dramatically \emph{worse} than the 8B model's 65.2\%
(Table~\ref{tab:model-size-scaling}).
The thinking tax for 27B ($75.7$\,pp relative to 8B non-thinking at the
same budget) is \textbf{2.6$\times$ larger} than for 8B ($28.8$\,pp).

The root cause is a scaling mismatch between reasoning chain length and
token budget.
The 27B model produces longer, more elaborate reasoning chains---a
consequence of its greater capacity and RLHF training.
At $b{=}512$, the 27B model's natural-stop rate is a mere \textbf{0.7\%}
(Table~\ref{tab:model-size-scaling}, footnote), meaning 99.3\% of its
responses are \emph{truncated projections}: the model ran out of tokens
mid-reasoning and was forced to output an answer.
By contrast, 8B achieves a 37.4\% natural-stop rate at the same budget.

This finding has important practical implications.
As reasoning models grow larger and more capable, their chains of thought
grow proportionally longer, requiring commensurately larger budgets to
avoid truncation.
The thinking tax does not diminish with scale---it \emph{worsens},
making adaptive budget allocation \emph{more} important, not less, for
frontier models.

\begin{figure}[t]
\centering
\includegraphics[width=0.72\textwidth]{fig_thinking_tax_model_size}
\caption{%
  \textbf{The thinking tax scales with model size.}
  Comparison of thinking-mode accuracy for Qwen3-8B (blue) and
  Qwen3.5-27B (red) at matched budgets on GSM8K.
  At $b{=}512$, 27B's accuracy (18.3\%) is 46.9\,pp below 8B's (65.2\%),
  because 27B's longer chains are more severely truncated.
  The 8B non-thinking baseline (green dashed) dominates both at all budgets.
}
\label{fig:thinking-tax-model-size}
\end{figure}

% ---------------------------------------------------------------------------
\subsection{Finding 5: 31.8\% of Questions Are Impossible at All Budgets}
\label{sec:finding:impossible}
% ---------------------------------------------------------------------------

To understand the efficiency ceiling of any budget-allocation scheme, we
categorize GSM8K problems by the \emph{minimum} budget at which Qwen3-8B
solves them correctly in thinking mode.

% Table 3: Thinking Efficiency Frontier
% Key result: 31.8% of questions are "impossible" at any tested budget
\begin{table}[t]
\centering
\caption{\textbf{Thinking Efficiency Frontier: difficulty distribution.}
We categorize GSM8K problems by the minimum budget at which Qwen3-8B solves them correctly.
31.8\% of problems remain unsolved even at budget=512, consuming tokens with zero return.
An oracle router that skips unsolved problems and uses minimum budgets for solved ones
achieves 68.2\% accuracy ($+$3.0pp over fixed-512) while using only 401 avg tokens (21.7\% savings).}
\label{tab:efficiency-frontier}
\small
\begin{tabular}{lrrrrr}
\toprule
\textbf{Category} & \textbf{Min Budget} & \textbf{Count} & \textbf{\%} & \textbf{Cumul. Acc} & \textbf{$\Delta$ Acc} \\
\midrule
Easy       & $\leq$128 & 155  & 11.8 & 11.8 & +11.8 \\
Medium     & $\leq$256 & 340  & 25.8 & 37.5 & +25.8 \\
Hard       & $\leq$512 & 405  & 30.7 & 68.2 & +30.7 \\
Impossible & ---       & 419  & 31.8 & ---  & --- \\
\midrule
\multicolumn{2}{l}{\textbf{Total}} & \textbf{1319} & 100 & & \\
\bottomrule
\end{tabular}
\end{table}


Table~\ref{tab:efficiency-frontier} and Figure~\ref{fig:difficulty-distribution}
reveal a four-tier difficulty distribution:
\begin{itemize}[nosep,leftmargin=*]
  \item \textbf{Easy} (11.8\%, $n{=}155$): solved at $b{\leq}128$.
    These require only a short chain.
  \item \textbf{Medium} (25.8\%, $n{=}340$): solved at $b{\leq}256$
    but not at 128.  A moderate reasoning chain suffices.
  \item \textbf{Hard} (30.7\%, $n{=}405$): solved only at $b{=}512$.
    These genuinely benefit from extended reasoning.
  \item \textbf{Impossible} (31.8\%, $n{=}419$): unsolved at \emph{any}
    tested budget ($b{\leq}512$).  Tokens spent on these problems
    yield zero return.
\end{itemize}

The ``impossible'' category is surprisingly large: nearly a third of
GSM8K is beyond the model's reach regardless of how much thinking budget
is allocated (up to 512 tokens).
An \emph{oracle router} that (i)~skips impossible questions entirely
and (ii)~assigns the minimum sufficient budget to solvable questions
achieves \textbf{68.2\%} accuracy ($+3.0$\,pp over fixed-512) while
consuming only \textbf{401} average tokens---a \textbf{21.7\%}
savings.\footnote{%
  The oracle accuracy exceeds fixed-512 because it avoids cases where
  extended reasoning causes the model to ``overthink'' and revise a
  correct intermediate answer.}

This analysis establishes both an opportunity and a ceiling for adaptive
methods.
The opportunity: a substantial fraction of compute is wasted on problems
that are either trivially easy (and need few tokens) or genuinely
impossible (and waste all tokens).
The ceiling: no budget controller operating within the $b{\leq}512$
regime can exceed 68.2\% accuracy on this benchmark, because 31.8\% of
problems are fundamentally beyond the model's capability.

\begin{figure}[t]
\centering
\includegraphics[width=0.72\textwidth]{fig3_difficulty_distribution}
\caption{%
  \textbf{Four-tier difficulty distribution of GSM8K problems} under Qwen3-8B thinking mode.
  Problems are categorized by the minimum budget at which they are solved correctly.
  31.8\% of problems are impossible at all tested budgets---tokens spent on these yield zero return.
  An ideal budget allocator would skip impossible questions and assign minimum budgets to the rest.
}
\label{fig:difficulty-distribution}
\end{figure}

\paragraph{Summary of findings.}
Taken together, our five findings paint a coherent picture:
\begin{enumerate}[nosep,leftmargin=*]
  \item The thinking tax is real and large (\S\ref{sec:finding:tax}):
    non-thinking mode dominates at all matched budgets.
  \item Models already signal when thinking was successful
    (\S\ref{sec:finding:natural-stop}): natural stop predicts 96.3\%
    accuracy at zero cost.
  \item High budgets are under-utilized (\S\ref{sec:finding:utilization}):
    up to 60\% of allocated tokens go unused.
  \item The tax worsens with scale (\S\ref{sec:finding:model-size}):
    larger models need proportionally larger budgets.
  \item A third of problems are impossible
    (\S\ref{sec:finding:impossible}): no budget helps.
\end{enumerate}
The natural question is: can we design a system that \emph{thinks only
when thinking actually helps}?
This is precisely the goal of \textsc{Town}, introduced next.

% ===========================================================================

Before introducing our method, we present a systematic empirical study of
the costs and failure modes of chain-of-thought reasoning under constrained
token budgets.  All experiments in this section use
\textbf{Qwen3-8B}~\cite{qwen3_2025} with its native thinking mode
(\texttt{enable\_thinking=True}) and \textbf{Qwen3.5-27B} for cross-scale
validation, evaluated on the full GSM8K test set ($n{=}1{,}319$)
unless otherwise noted.\footnote{%
  A 200-sample subset (seed 42) is used for the budget sweep in
  Table~\ref{tab:thinking-tax}; all other analyses use the complete set.}
We control test-time compute via the \texttt{max\_new\_tokens} parameter,
which caps the \emph{total} number of newly generated tokens---including
both the thinking trace (\texttt{<think>...<\!/think>}) and the final answer.
Crucially, thinking mode and non-thinking mode operate under the
\emph{same} total output-token budget~$b$, ensuring a fair comparison:
\texttt{think@$b$} must fit its entire reasoning chain \emph{plus}
its answer within $b$~tokens, while \texttt{nothink@$b$} uses all
$b$~tokens for the answer directly.
Non-thinking results are obtained by setting
\texttt{enable\_thinking=False} in the Qwen3 chat template, which
suppresses the thinking block entirely.
Greedy decoding ($\tau{=}0$) is used throughout.

Our five findings collectively establish that (i)~thinking mode imposes a
severe \emph{tax} at practical token budgets, (ii)~models already emit
reliable signals indicating when this tax is unnecessary, and
(iii)~a large fraction of compute is wasted on questions the model
cannot solve regardless of budget.
These findings motivate the design of \textsc{Town} in
\S\ref{sec:method}.

% ---------------------------------------------------------------------------
\subsection{Finding 1: Non-Thinking Beats Thinking at All Matched Budgets}
\label{sec:finding:tax}
% ---------------------------------------------------------------------------

Our most striking result is that, when the output-token budget is held
constant, \emph{non-thinking mode uniformly dominates thinking mode}
across every budget level we tested.
Table~\ref{tab:thinking-tax} summarizes the comparison.

% Table 1: The Thinking Tax — Comprehensive comparison at matched budgets
\begin{table}[t]
\centering
\caption{\textbf{The Thinking Tax: Non-thinking mode dramatically outperforms thinking mode at matched token budgets.}
Results on GSM8K with Qwen3-8B.
Rows without a marker use the 200-sample subset (seed=42); starred ($\star$) rows use the full test set ($n{=}1{,}319$).
Qwen3.5-27B results use the full set.
\textbf{At every budget, non-thinking achieves higher accuracy while using fewer tokens.}}
\label{tab:thinking-tax}
\small
\begin{tabular}{ll rrrr}
\toprule
\textbf{Model} & \textbf{Config} & \textbf{Budget} & \textbf{Accuracy} & \textbf{Avg Tok} & \textbf{Early Stop} \\
\midrule
\multirow{10}{*}{\textbf{Qwen3-8B}}
& \texttt{nothink@128}      & 128 & 54.5\%  & 111 & 43.5\% \\
& \texttt{nothink@128$\star$} & 128 & 50.8\%  & 113 & --- \\
& \texttt{nothink@256}      & 256 & 89.0\%  & 140 & 92.0\% \\
& \texttt{nothink@256$\star$} & 256 & \textbf{87.5\%}  & \textbf{146} & 88.8\% \\
& \texttt{nothink@512}      & 512 & 94.0\%  & 145 & 99.5\% \\
\cmidrule{2-6}
& \texttt{think@128}        & 128 & 2.0\%   & 128 & 0.0\% \\
& \texttt{think@128$\star$}  & 128 & 3.0\%   & 128 & --- \\
& \texttt{think@256}        & 256 & 22.0\%  & 255 & 2.0\% \\
& \texttt{think@256$\star$}  & 256 & 18.0\%  & 255 & --- \\
& \texttt{think@512}        & 512 & 66.5\%  & 442 & 47.5\% \\
& \texttt{think@512$\star$}  & 512 & 65.2\%  & 460 & 37.4\% \\
\midrule
\multirow{3}{*}{\textbf{Qwen3.5-27B$\star$}}
& \texttt{think@128}        & 128 & 3.6\%   & 144 & 0.0\% \\
& \texttt{think@256}        & 256 & 7.9\%   & 272 & 0.0\% \\
& \texttt{think@512}        & 512 & 18.3\%  & 528 & 0.7\% \\
\bottomrule
\end{tabular}
\vspace{-2mm}
\end{table}


At budget $b{=}128$, non-thinking achieves \textbf{54.5\%} accuracy while
thinking yields only \textbf{2.0\%}---a 52.5 percentage-point (pp) gap.
At $b{=}256$, the gap narrows but remains enormous:
\textbf{89.0\%} vs.\ \textbf{22.0\%} ($+67.0$\,pp).
Even at $b{=}512$, where thinking has enough room for short chains,
non-thinking still leads by a large margin:
\textbf{94.0\%} vs.\ \textbf{66.5\%} ($+27.5$\,pp).
These results are confirmed on the full GSM8K test set ($n{=}1{,}319$,
$\star$ rows in Table~\ref{tab:thinking-tax}):
\texttt{nothink@256$\star$} achieves \textbf{87.5\%} vs.\ only
\textbf{18.0\%} for \texttt{think@256$\star$}, a
\textbf{69.5\,pp} gap at matched budget.

The root cause is straightforward: a thinking model must first produce a
reasoning chain \emph{inside} its $b$-token budget before emitting an answer.
At low budgets, the chain consumes nearly all available tokens, leaving
insufficient room for a well-formed answer.
A non-thinking model, by contrast, directly generates the solution,
allocating its full budget to answer-bearing content.
This difference manifests in the \emph{Early Stop} column of
Table~\ref{tab:thinking-tax}: at $b{=}256$, 92.0\% of non-thinking
samples terminate naturally before the budget, whereas only 2.0\% of
thinking samples do.
In effect, the non-thinking model has already ``solved'' the problem and
moved on, while the thinking model is still mid-chain.

% Figure 1 (nothink vs thinking) is presented in the Introduction.
% See Figure~\ref{fig:nothink-vs-thinking} for the visual comparison.

We emphasize that this result does \emph{not} imply that chain-of-thought
reasoning is useless---at sufficiently large budgets (e.g.,\ $b{=}4096$),
thinking mode recovers and eventually surpasses non-thinking.
Rather, it shows that under the budget-constrained regime that is
increasingly relevant for cost-sensitive deployment, the overhead of
producing a reasoning chain is a \emph{tax} that must be justified by
the problem's difficulty.

% ---------------------------------------------------------------------------
\subsection{Finding 2: Natural Stop Is a Free Confidence Oracle}
\label{sec:finding:natural-stop}
% ---------------------------------------------------------------------------

When a thinking model generates a complete reasoning chain and emits
a ``Final answer'' marker \emph{before} exhausting its token budget,
we say the sample exhibits a \emph{natural stop}.
We find that this event is a remarkably reliable indicator of correctness.

% Table 1: Natural Stop as Confidence Oracle
% Key result: natural-stop samples achieve 99.0% accuracy at budget=512 (HF engine)
% Updated 2026-04-09: switched from vLLM+projection to HF engine for consistency
\begin{table}[t]
\centering
\caption{\textbf{Natural stop as a free confidence oracle.}
When a thinking model naturally terminates before exhausting its token budget,
the resulting answer is correct with remarkably high probability.
This phenomenon is consistent across models and benchmarks.
\emph{NatStop}: fraction of samples that stop early ($<$95\% of budget used).
\emph{Acc$_\text{nat}$}: accuracy among early-stop samples.
\emph{Acc$_\text{hit}$}: accuracy among budget-exhausting samples.
All Qwen3-8B results use HF engine with greedy decoding.
DeepSeek-R1 Acc splits are omitted due to small sample sizes ($n{=}40$--$80$).}
\label{tab:natural-stop-oracle}
\small
\setlength{\tabcolsep}{3pt}
\begin{tabular}{@{}lrrrrrr@{}}
\toprule
\textbf{Model / Benchmark} & \textbf{Bud.} & \textbf{Acc} & \textbf{NatStop} & \textbf{Acc$_\text{nat}$} & \textbf{Acc$_\text{hit}$} & \textbf{Gap} \\
\midrule
\multirow{3}{*}{Qwen3-8B / GSM8K (n=1319)}
 & 128  &  3.0 &  0.0 & --- &  3.0 & --- \\
 & 256  & 18.0 &  1.4 & 100.0 & 16.8 & +83.2 \\
 & 512  & 56.9 & 37.4 & 99.0 & 31.8 & \textbf{+67.2} \\
\midrule
\multirow{3}{*}{DeepSeek-R1-8B / GSM8K (n=40)}
 & 256  & 22.5 & 15.0 & --- & --- & --- \\
 & 512  & 55.0 & 85.0 & --- & --- & --- \\
 & 1024 & 65.0 & 97.5 & --- & --- & --- \\
\midrule
\multirow{3}{*}{DeepSeek-R1-8B / MATH500 (n=80)}
 & 1024 & 25.0 & 31.2 & --- & --- & --- \\
 & 2048 & 38.7 & 55.0 & --- & --- & --- \\
 & 4096 & 46.2 & 76.2 & --- & --- & --- \\
\bottomrule
\end{tabular}
\end{table}


At $b{=}512$ with Qwen3-8B, 37.4\% of samples terminate before consuming
95\% of their token budget---we call these \emph{natural stops}
(Table~\ref{tab:natural-stop-oracle}).
Among these 475 samples, accuracy is \textbf{96.3\%}---nearly perfect.
Among samples that exhaust the budget (i.e., are truncated), accuracy is
only 47.5\%, yielding an accuracy gap of \textbf{$+$48.8\,pp}
(Figure~\ref{fig:natural-stop-split}).
A stricter definition---requiring an explicit ``Final answer'' marker
(\texttt{has\_final})---selects only 6.1\% of samples but achieves
93.8\% accuracy, further confirming the signal's reliability.

This phenomenon admits a simple explanation.
A natural stop means the model's \emph{own} generation process
determined that the reasoning chain was complete---the model ``believes''
it has reached a satisfactory answer.
While LLM confidence is notoriously miscalibrated in
general~\cite{chen2024overthinking}, the natural-stop signal is
structurally different from soft confidence scores: it is a
\emph{binary}, \emph{endogenous} event that does not require access to
logits, hidden states, or any post-hoc calibration.
It is, in a precise sense, a \emph{free} confidence oracle---one that
the model provides at zero additional inference cost.

\begin{remark}[Generality]
The natural-stop oracle effect is not limited to Qwen3-8B.
Table~\ref{tab:natural-stop-oracle} shows that DeepSeek-R1-8B exhibits
similar behavior: at $b{=}512$, 85.0\% of GSM8K samples stop naturally,
and on the harder MATH500 benchmark, natural-stop rates increase steadily
with budget (31.2\% $\to$ 55.0\% $\to$ 76.2\% as $b$ scales from 1024
to 4096).
This universality suggests that the phenomenon is an intrinsic property
of reasoning-trained models, not an artifact of a specific architecture.
\end{remark}

\begin{figure}[t]
\centering
\includegraphics[width=0.72\textwidth]{fig9_acc_split_natural_vs_hit}
\caption{%
  \textbf{Natural stop is a free confidence oracle.}
  Accuracy decomposition by stopping behavior for Qwen3-8B on GSM8K.
  Samples that stop naturally before the budget (blue) achieve dramatically
  higher accuracy than those that exhaust the budget (red), across all
  budget levels.  At $b{=}512$, natural-stop samples reach 96.3\% accuracy
  vs.\ 47.5\% for truncated samples---a 48.8\,pp gap.
}
\label{fig:natural-stop-split}
\end{figure}

% ---------------------------------------------------------------------------
\subsection{Finding 3: Token Utilization Collapses at High Budgets}
\label{sec:finding:utilization}
% ---------------------------------------------------------------------------

A complementary perspective on the natural-stop phenomenon comes from
examining \emph{token utilization}: the fraction of the allocated budget
that the model actually consumes.

% Table 2: Token Utilization Decreases at Higher Budgets
% Key result: utilization drops from 100% to 38% as budget increases
\begin{table}[t]
\centering
\caption{\textbf{Token utilization ratio decreases with budget.}
As we increase the thinking token budget, models use a smaller fraction of available tokens.
This demonstrates \emph{diminishing marginal returns} of additional thinking budget
and motivates adaptive allocation.
Utilization = avg tokens used / budget.
$^{\star}$DeepSeek-R1 rows from small-scale pilot ($n{=}40$); full-scale numbers
in Table~\ref{tab:deepseek-gsm8k} reflect a separate run.}
\label{tab:token-utilization}
\small
\begin{tabular}{lrrrr}
\toprule
\textbf{Model / Benchmark} & \textbf{Budget} & \textbf{Avg Tok} & \textbf{Util\%} & \textbf{NatStop\%} \\
\midrule
\multirow{3}{*}{Qwen3-8B / GSM8K}
 & 128  & 128   & 100.0 &  0.0 \\
 & 256  & 255   &  99.7 &  1.4 \\
 & 512  & 460   &  89.8 & 37.4 \\
\midrule
\multirow{3}{*}{DeepSeek-R1 / GSM8K$^{\star}$}
 & 256  & 264   & 103.0 & 15.0 \\
 & 512  & 376   &  73.5 & 85.0 \\
 & 1024 & 393   &  \textbf{38.3} & 97.5 \\
\midrule
\multirow{3}{*}{DeepSeek-R1 / MATH500}
 & 1024 & 910   &  88.8 & 31.2 \\
 & 2048 & 1551  &  75.7 & 55.0 \\
 & 4096 & 2353  &  57.4 & 76.2 \\
\bottomrule
\end{tabular}
\vspace{-2mm}
\end{table}


Table~\ref{tab:token-utilization} reveals a striking pattern.
At $b{=}128$, Qwen3-8B uses 100\% of its budget---every token is
consumed, and the model is almost certainly still mid-reasoning when
generation is forcibly terminated.
At $b{=}256$, utilization remains near-perfect (99.7\%), with only 1.4\%
of samples stopping naturally.
But at $b{=}512$, utilization drops to \textbf{89.8\%}, with an average
of only 460 tokens consumed out of 512 available.
The model is already ``paying'' for $\sim$52 tokens it does not need.

The pattern is even more dramatic for DeepSeek-R1 on GSM8K: utilization
plummets from 103.0\% at $b{=}256$ (the model slightly exceeds its
budget due to token-counting granularity) to \textbf{38.3\%} at
$b{=}1{,}024$, with 97.5\% of samples stopping naturally.
This means the model uses fewer than 400 tokens on average when allocated
1{,}024---more than 60\% of the budget is wasted.

This utilization collapse is a direct consequence of the heterogeneous
difficulty distribution in typical benchmarks: many problems are easy
enough that a short reasoning chain suffices, and additional budget
yields no benefit.
It provides strong empirical motivation for adaptive budget allocation:
rather than setting a single high budget and accepting low utilization,
one should match the budget to the problem's actual difficulty.

% ---------------------------------------------------------------------------
\subsection{Finding 4: The Thinking Tax Scales with Model Size}
\label{sec:finding:model-size}
% ---------------------------------------------------------------------------

One might hope that the thinking tax diminishes with more capable models.
We test this hypothesis by comparing Qwen3-8B and Qwen3.5-27B in
thinking mode at identical budgets.
The results are counterintuitive.

% Table: Thinking Tax scales with model size (8B, 9B, 27B)
% Updated 2026-04-08 with gap-fill results: 8B nothink@512/1024/2048, 9B think@512/nothink@256/512, 8B think@1024
\begin{table}[t]
\centering
\caption{\textbf{The Thinking Tax worsens with model size} on GSM8K ($n{=}1{,}319$).
At budget=512, thinking-mode accuracy collapses from 56.9\% (8B) to 18.3\% (27B),
while non-thinking accuracy remains high across all scales:
8B=93.1\%, 9B=93.2\%, 27B=\textbf{95.5\%}.
The 9B model exhibits the largest tax (\textbf{77.7\,pp}), confirming that
inverse scaling is driven by chain-length growth, not capability loss.}
\label{tab:model-size-scaling}
\small
\begin{tabular}{l rr rr rr}
\toprule
 & \multicolumn{2}{c}{\textbf{8B}} & \multicolumn{2}{c}{\textbf{9B}} & \multicolumn{2}{c}{\textbf{27B}} \\
\cmidrule(lr){2-3} \cmidrule(lr){4-5} \cmidrule(lr){6-7}
\textbf{Budget} & \textbf{Think} & \textbf{NoThink} & \textbf{Think} & \textbf{NoThink} & \textbf{Think} & \textbf{NoThink} \\
\midrule
128 & 3.0\%  & 50.8\% & ---    & ---    & 3.6\%  & 9.9\% \\
256 & 18.0\% & 87.5\% & 3.4\%  & 61.2\% & 7.9\%  & 65.1\% \\
512 & 56.9\% & 93.1\% & 15.5\% & 93.2\% & 18.3\% & \textbf{95.5\%} \\
1024 & 86.1\% & 93.1\% & 41.8\% & ---$^\ddagger$ & ${\sim}$49\%$^\dagger$ & --- \\
2048 & ---$^\ddagger$ & 93.1\% & 66.8\% & ---    & ---    & --- \\
\midrule
\multicolumn{7}{l}{\textit{Thinking Tax at budget=512 (NoThink@512 $-$ Think@512):}} \\
\midrule
Tax & \multicolumn{6}{c}{8B: \textbf{36.2\,pp} (93.1\%$-$56.9\%); \quad 9B: \textbf{77.7\,pp} (93.2\%$-$15.5\%); \quad 27B: \textbf{77.2\,pp} (95.5\%$-$18.3\%)} \\
\bottomrule
\end{tabular}
\vspace{1mm}

\footnotesize
$^\dagger$Checkpoint at $n{=}200$; full-dataset experiment pending.\\
$^\dagger$Checkpoint at $n{=}200$; full-dataset experiment pending.
$^\ddagger$8B think@2048 in progress.\\
\textit{Root cause}: natural-stop rates at budget=512: 8B=37.4\%, 9B${\approx}$0.5\%, 27B=0.7\%.\\
At $b{=}1024$: 8B think=86.1\% vs nothink=93.1\% (gap 7.0\,pp); crossover not yet reached.\\
Nothink saturates: 8B nothink@512/1024/2048 = 93.1\%/93.1\%/93.1\% (avg ${\sim}$152 tokens).\\
\vspace{-3mm}
\end{table}


At $b{=}512$, the 27B model achieves only \textbf{18.3\%} accuracy in
thinking mode---dramatically \emph{worse} than the 8B model's 65.2\%
(Table~\ref{tab:model-size-scaling}).
The thinking tax for 27B ($75.7$\,pp relative to 8B non-thinking at the
same budget) is \textbf{2.6$\times$ larger} than for 8B ($28.8$\,pp).

The root cause is a scaling mismatch between reasoning chain length and
token budget.
The 27B model produces longer, more elaborate reasoning chains---a
consequence of its greater capacity and RLHF training.
At $b{=}512$, the 27B model's natural-stop rate is a mere \textbf{0.7\%}
(Table~\ref{tab:model-size-scaling}, footnote), meaning 99.3\% of its
responses are \emph{truncated projections}: the model ran out of tokens
mid-reasoning and was forced to output an answer.
By contrast, 8B achieves a 37.4\% natural-stop rate at the same budget.

This finding has important practical implications.
As reasoning models grow larger and more capable, their chains of thought
grow proportionally longer, requiring commensurately larger budgets to
avoid truncation.
The thinking tax does not diminish with scale---it \emph{worsens},
making adaptive budget allocation \emph{more} important, not less, for
frontier models.

\begin{figure}[t]
\centering
\includegraphics[width=0.72\textwidth]{fig_thinking_tax_model_size}
\caption{%
  \textbf{The thinking tax scales with model size.}
  Comparison of thinking-mode accuracy for Qwen3-8B (blue) and
  Qwen3.5-27B (red) at matched budgets on GSM8K.
  At $b{=}512$, 27B's accuracy (18.3\%) is 46.9\,pp below 8B's (65.2\%),
  because 27B's longer chains are more severely truncated.
  The 8B non-thinking baseline (green dashed) dominates both at all budgets.
}
\label{fig:thinking-tax-model-size}
\end{figure}

% ---------------------------------------------------------------------------
\subsection{Finding 5: 31.8\% of Questions Are Impossible at All Budgets}
\label{sec:finding:impossible}
% ---------------------------------------------------------------------------

To understand the efficiency ceiling of any budget-allocation scheme, we
categorize GSM8K problems by the \emph{minimum} budget at which Qwen3-8B
solves them correctly in thinking mode.

% Table 3: Thinking Efficiency Frontier
% Key result: 31.8% of questions are "impossible" at any tested budget
\begin{table}[t]
\centering
\caption{\textbf{Thinking Efficiency Frontier: difficulty distribution.}
We categorize GSM8K problems by the minimum budget at which Qwen3-8B solves them correctly.
31.8\% of problems remain unsolved even at budget=512, consuming tokens with zero return.
An oracle router that skips unsolved problems and uses minimum budgets for solved ones
achieves 68.2\% accuracy ($+$3.0pp over fixed-512) while using only 401 avg tokens (21.7\% savings).}
\label{tab:efficiency-frontier}
\small
\begin{tabular}{lrrrrr}
\toprule
\textbf{Category} & \textbf{Min Budget} & \textbf{Count} & \textbf{\%} & \textbf{Cumul. Acc} & \textbf{$\Delta$ Acc} \\
\midrule
Easy       & $\leq$128 & 155  & 11.8 & 11.8 & +11.8 \\
Medium     & $\leq$256 & 340  & 25.8 & 37.5 & +25.8 \\
Hard       & $\leq$512 & 405  & 30.7 & 68.2 & +30.7 \\
Impossible & ---       & 419  & 31.8 & ---  & --- \\
\midrule
\multicolumn{2}{l}{\textbf{Total}} & \textbf{1319} & 100 & & \\
\bottomrule
\end{tabular}
\end{table}


Table~\ref{tab:efficiency-frontier} and Figure~\ref{fig:difficulty-distribution}
reveal a four-tier difficulty distribution:
\begin{itemize}[nosep,leftmargin=*]
  \item \textbf{Easy} (11.8\%, $n{=}155$): solved at $b{\leq}128$.
    These require only a short chain.
  \item \textbf{Medium} (25.8\%, $n{=}340$): solved at $b{\leq}256$
    but not at 128.  A moderate reasoning chain suffices.
  \item \textbf{Hard} (30.7\%, $n{=}405$): solved only at $b{=}512$.
    These genuinely benefit from extended reasoning.
  \item \textbf{Impossible} (31.8\%, $n{=}419$): unsolved at \emph{any}
    tested budget ($b{\leq}512$).  Tokens spent on these problems
    yield zero return.
\end{itemize}

The ``impossible'' category is surprisingly large: nearly a third of
GSM8K is beyond the model's reach regardless of how much thinking budget
is allocated (up to 512 tokens).
An \emph{oracle router} that (i)~skips impossible questions entirely
and (ii)~assigns the minimum sufficient budget to solvable questions
achieves \textbf{68.2\%} accuracy ($+3.0$\,pp over fixed-512) while
consuming only \textbf{401} average tokens---a \textbf{21.7\%}
savings.\footnote{%
  The oracle accuracy exceeds fixed-512 because it avoids cases where
  extended reasoning causes the model to ``overthink'' and revise a
  correct intermediate answer.}

This analysis establishes both an opportunity and a ceiling for adaptive
methods.
The opportunity: a substantial fraction of compute is wasted on problems
that are either trivially easy (and need few tokens) or genuinely
impossible (and waste all tokens).
The ceiling: no budget controller operating within the $b{\leq}512$
regime can exceed 68.2\% accuracy on this benchmark, because 31.8\% of
problems are fundamentally beyond the model's capability.

\begin{figure}[t]
\centering
\includegraphics[width=0.72\textwidth]{fig3_difficulty_distribution}
\caption{%
  \textbf{Four-tier difficulty distribution of GSM8K problems} under Qwen3-8B thinking mode.
  Problems are categorized by the minimum budget at which they are solved correctly.
  31.8\% of problems are impossible at all tested budgets---tokens spent on these yield zero return.
  An ideal budget allocator would skip impossible questions and assign minimum budgets to the rest.
}
\label{fig:difficulty-distribution}
\end{figure}

\paragraph{Summary of findings.}
Taken together, our five findings paint a coherent picture:
\begin{enumerate}[nosep,leftmargin=*]
  \item The thinking tax is real and large (\S\ref{sec:finding:tax}):
    non-thinking mode dominates at all matched budgets.
  \item Models already signal when thinking was successful
    (\S\ref{sec:finding:natural-stop}): natural stop predicts 96.3\%
    accuracy at zero cost.
  \item High budgets are under-utilized (\S\ref{sec:finding:utilization}):
    up to 60\% of allocated tokens go unused.
  \item The tax worsens with scale (\S\ref{sec:finding:model-size}):
    larger models need proportionally larger budgets.
  \item A third of problems are impossible
    (\S\ref{sec:finding:impossible}): no budget helps.
\end{enumerate}
The natural question is: can we design a system that \emph{thinks only
when thinking actually helps}?
This is precisely the goal of \textsc{Town}, introduced next.
